\documentclass[11pt,a4paper]{article}
\usepackage[margin=0.85in,top=1in,bottom=1in]{geometry}
\usepackage{amsmath,amssymb,amsfonts}
\usepackage{graphicx}
\usepackage{float}
\usepackage{enumitem}
\usepackage{microtype}
\usepackage{caption}
\usepackage[hidelinks]{hyperref}
\usepackage[most,breakable]{tcolorbox}
\tcbuselibrary{breakable,skins}
\usepackage[utf8]{inputenc}
\usepackage[T1]{fontenc}
\usepackage{lmodern}
\captionsetup{font=small,labelfont=bf,skip=4pt}
\setlist{leftmargin=1.5em,topsep=2pt}

%% Breakable card — prevents content cutoff across pages
\newtcolorbox{solcard}[3]{
  breakable, enhanced,
  colback=white, colframe=gray!50, arc=2pt,
  title={\textbf{#1}\hfill\textsf{\small #3\,pts}},
  fonttitle=\normalsize\sffamily\bfseries,
  before upper={\small\textit{#2}\par\smallskip\hrule height 0.4pt\smallskip},
  top=4pt, bottom=6pt, left=7pt, right=7pt,
  parbox=false,
}

\title{\textbf{Propagation of detonation waves --- Solution}}
\author{\normalsize Y23 (2023) — English translation}
\date{}
\begin{document}\maketitle\vspace{-0.5em}
\begin{solcard}{A1}{Write down a system of laws of conservation of mass and momentum for the volume of gas passing through the wave front of the shock wave. Determine from it the speed of movement of the wave front $c$.}{1.00}

Let us move to a reference frame that moves with the velocity of the shock wave front $c$. In this CO, a gas with pressure $p_0$ and density $\rho_0$ impinges on a stationary front, and after the front moves with a certain speed $u$. Throughout the problem we will use this reference system. In time $dt$, mass $\rho_0  c dt$ passes through a unit area of the flow, and mass $\rho_1 u dt$ exits, therefore the law of conservation of mass $$ \rho_0 c = \rho_1 u. $$Momentum one can obtain, multiply by mass on correspond velocity. At the same time the momentum of the gas varies as the wave front passes due to the pressure difference by the amount $(p_0 - p_1) dt$, therefore $$ (\rho_1 u^2 - \rho_0 c^2 )dt = (p_0 - p_1 )dt, $$and finally $$ p_0 + \rho_0 c^2 = p_1 + \rho_1 u^2. $$Eliminate velocity $u$ with means law conservation mass: $$ u = c \frac{\rho_0}{\rho_1} $$and obtain $$ p_0 + \rho_0 c^2 = p_1 + \frac{\rho_0^2}{\rho_1} c^2, $$from which velocity wave $$ c = \sqrt{\frac{p_1 - p_0}{\rho_1 - \rho_0} \frac{\rho_1}{\rho_0}}. $$

\par\smallskip\noindent\textbf{Answer:} $$ \rho_0 c = \rho_1 u, \quad p_0 + \rho_0 c^2 = p_1 + \frac{\rho_0^2}{\rho_1} c^2, \quad c = \sqrt{\frac{p_1 - p_0}{\rho_1 - \rho_0} \frac{\rho_1}{\rho_0}}. $$
\end{solcard}

\begin{solcard}{A2}{Write down the law of conservation of energy for an element of the volume of gas passing through the wave front. Obtain from it an additional equation relating the parameters $p_0$, $\rho_0$, $p_1$, $\rho_1$, the speed of movement of the wave front $c$ and the adiabatic exponent $\gamma$.}{1.00}

Let us consider the mass $dm$ of gas crossing the boundary. Its kinetic energy changes by $$ \Delta E_{\text{to}} = \left( \frac{u^2}{2}  -\frac{c^2}{2}\right) dm, $$internal energy on $$ \Delta U =  C_V \frac{dm}{\mu} (T_1 - T_0 )= C_V \frac{dm}{R} \left( \frac{p_1}{\rho_1} - \frac{p_0}{\rho_0}\right). $$Here $C_V$ is the molar heat capacity of the gas at constant volume, and temperature express through pressure with means equation state ideal gas. Above gas perform work $$ A =p_0 \frac{dm}{\rho_0} -p_1 \frac{dm}{\rho_1}. $$Therefore energy conservation law \begin{align*} A \&= \Delta E_{\text{to}} + \Delta U, \\ p_0 \frac{dm}{\rho_0} -p_1 \frac{dm}{\rho_1}\& = \left( \frac{u^2}{2}  -\frac{c^2}{2}\right) dm + C_V \frac{dm}{R} \left( \frac{p_1}{\rho_1} - \frac{p_0}{\rho_0}\right). \end{align*} We cancel by the mass and transfer the remaining terms relating to the gas through the wave front to one side: $$ \frac{p_0}{\rho_0} + \frac{C_V}{R} \frac{p_0}{\rho_0} + \frac{c^2}{2} = \frac{p_1}{\rho_1} + \frac{C_V}{R} \frac{p_1}{\rho_1} + \frac{u^2}{2}. $$This relation is Bernoulli's equation, in which internal energy gas express through pressure and density. Note, that $$ \frac{C_V}{R} + 1 = \frac{C_P}{R} = \frac{C_P}{C_P-C_V} = \frac{\gamma}{\gamma - 1}, $$and also substitute expression for velocity $u$ through velocity wave front $c$ and we get

\par\smallskip\noindent\textbf{Answer:} $$ \frac{\gamma}{\gamma - 1} \frac{p_0}{\rho_0} + \frac{c^2}{2} = \frac{\gamma}{\gamma - 1} \frac{p_1}{\rho_1} + \frac{\rho_0^2}{\rho_1^2}\frac{c^2}{2} $$
\end{solcard}

\begin{solcard}{A3}{Combining the results obtained in steps A1 and A2, express the pressure ratio $p_1/p_0$ in terms of the adiabatic exponent $\gamma$ and the compaction coefficient $k=\rho_1/\rho_0$. This relationship is called the adiabatic equation for a shock wave.}{0.50}

Let us substitute the expression for the wave front speed into the relation from the previous paragraph: $$ \frac{\gamma}{\gamma - 1}\left( \frac{p_1}{\rho_1}- \frac{p_0}{\rho_0}\right) = \frac{1}{2}\left(1-\frac{\rho_0^2}{\rho_1^2} \right) \frac{p_1- p_0}{\rho_1 - \rho_0}\frac{\rho_1}{\rho_0}. $$Simplify right part $$ \frac{\gamma}{\gamma - 1}\left( \frac{p_1}{\rho_1}- \frac{p_0}{\rho_0}\right) = \frac{\rho_1 + \rho_0}{2\rho_0 \rho_1} (p_1 - p_0) , $$group the terms with $p_0$ and with $p_1$: $$ p_1\left(\frac{\gamma}{\gamma - 1} \frac{1}{\rho_1} -\frac{\rho_1 + \rho_0}{2 \rho_0 \rho_1} \right) = p_0\left(\frac{\gamma}{\gamma - 1} \frac{1}{\rho_0} - \frac{\rho_1 + \rho_0}{2 \rho_0 \rho_1} \right), $$from which the required ratio $$ \frac{p_1}{p_0}= \frac{\frac{\gamma}{\gamma - 1} \rho_1 - \frac{1}{2}(\rho_0 + \rho_1)}{\frac{\gamma}{\gamma - 1} \rho_0 - \frac{1}{2}(\rho_0 + \rho_1)} = \frac{2 \gamma k - (\gamma -1) (1 + k)}{2 \gamma - (\gamma - 1)(1 + k) }. $$

\par\smallskip\noindent\textbf{Answer:} $$ \frac{p_1}{p_0} = \frac{(\gamma + 1)k  - (\gamma -1)}{\gamma + 1 - (\gamma - 1)k }. $$
\end{solcard}

\begin{solcard}{A4}{Determine the maximum possible value of the compaction coefficient $k_{max}$. Express your answer using $\gamma$. Calculate the answer for an ideal diatomic gas.}{0.50}

From the formula for the previous paragraph we see that the pressure $p_1$ goes to infinity at $k = \dfrac{\gamma + 1}{\gamma -1}$. This is the maximum value $k$. For diatomic gas $\gamma = 7/5$, $k_{max} = 6$.

\par\smallskip\noindent\textbf{Answer:} $$ k_{max} = \frac{\gamma + 1}{\gamma -1} = 6. $$
\end{solcard}

\begin{solcard}{B1}{Write down the law of conservation of energy for a detonation wave. Obtain from it an equation relating the parameters $p_0$, $\rho_0$, $p_1$, $\rho_1$, the speed of movement of the wave front $c$, the specific thermal effect of the reaction $q$, as well as the adiabatic exponents $\gamma_0$ and $\gamma_1$.}{0.60}

The law of conservation of energy is similar to point $\mathrm{B3}$, but you need to add the released heat $q dm$ to the work of pressure forces and take into account that the molar heat capacities of the original gas and reaction products differ: $$ p_0 \frac{dm}{\rho_0} -p_1 \frac{dm}{\rho_1} + q dm = \left( \frac{u^2}{2}  -\frac{c^2}{2}\right) dm + C_{V1} \frac{dm}{R}  \frac{p_1}{\rho_1} - C_{V0} \frac{dm}{R} \frac{p_0}{\rho_0}. $$Express heat capacity gas through exponent/index adiabat, velocity gas $u$ through the wave front velocity and rearrange the terms similarly to part A:

\par\smallskip\noindent\textbf{Answer:} $$ \frac{\gamma_0}{\gamma_0 - 1} \frac{p_0}{\rho_0} + \frac{c^2}{2} + q = \frac{\gamma_1}{\gamma_1 -1} \frac{p_1}{\rho_1} + \frac{c^2}{2} \frac{\rho_0^2}{\rho_1^2}. $$
\end{solcard}

\begin{solcard}{B2}{The explicit dependence $p_1(v_1)_\text{det}$ for a strong detonation wave can be represented as: $$p_1=\cfrac{q}{v_0}\cdot A $$Determine $A$. Answer express through $\gamma_1$ and ratio $v_1/v_0$. This relationship is called the detonation adiabatic equation.}{0.60}

In the strong wave approximation, the wavefront speed is $$ c^2 = \frac{p_1}{\rho_1 - \rho_0} \frac{\rho_1}{\rho_0} = p_1 \frac{v_0^2}{v_0 - v_1}. $$Substitute it in energy conservation law, neglect at the same time term $\frac{p_0}{\rho_0}$: $$ q + p_1 \frac{v_0^2}{2(v_0 - v_1)} = \frac{\gamma_1}{\gamma_1 -1}p_1 v_1 + \frac{v_1^2 }{v_0^2} p_1 \frac{v_0^2}{2(v_0 - v_1)} . $$Hence $$ p_1 \left( \frac{\gamma_1}{\gamma_1 -1} v_1  + \frac{v_1^2 - v_0^2}{2 (v_0 - v_1)}\right) = q, $$$$ p_1 = \frac{q}{\frac{\gamma_1}{\gamma_1 -1} v_1  -  \frac{1}{2} (v_0 + v_1)} = \frac{q}{v_0} \frac{2}{\frac{\gamma_1 + 1}{\gamma_1 -1} \frac{v_1}{v_0} - 1} $$

\par\smallskip\noindent\textbf{Answer:} $$ A =  \frac{2}{\frac{\gamma_1 + 1}{\gamma_1 -1} \frac{v_1}{v_0} - 1}. $$
\end{solcard}

\begin{solcard}{B3}{On one graph, qualitatively depict the dependence $p_1(v_1)_\text{det}$, as well as the dependence $p_1(v_1)_\text{imp}$ for different values of the speed of movement of the wave front $c$.}{0.60}

The graph of the dependence $p_1(v_1)_{\text{det}}$ is a hyperbola whose asymptote is at $v_1 = v_0 \dfrac{\gamma_1 -1}{\gamma_1 + 1}$. The dependence of pressure on specific volume, expressed in terms of the speed of sound, has the form $$ p_1(v_1)_{\text{imp}} = \frac{c^2}{v_0} \left( 1 - \frac{v_1}{v_0}\right), $$its graph is a linear function, vanishing at $v_1 = v_0$.

\begin{figure}[H]
  \centering
  \includegraphics[width=0.72\textwidth]{T4_S_EN_assets/fig1.jpg}
\end{figure}
\end{solcard}

\begin{solcard}{B4}{Determine the parameters $p_1$ and $v_1$ of the Chapman-Jouguet point. Express your answer using $v_0$, $\gamma_1$, $q$.}{1.40}

Let us equate the expressions for pressure through the wave speed and from the shock adiabatic equation. $$ \frac{c^2}{v_0} \left( 1 - \frac{v_1}{v_0}\right)= \frac{2q}{v_0} \frac{1}{\frac{\gamma_1 +1}{\gamma_1 - 1} \frac{v_1}{v_0} - 1}, $$obtain square equation on ratio $v_1/v_0$: $$ \left(1- \frac{v_1}{v_0} \right)\left( \frac{\gamma_1 +1}{\gamma_1 - 1} \frac{v_1}{v_0} - 1\right) = \frac{2q}{c^2}. $$This equation be have unique solution, if right part equal maximum value quadratic function in left part. This maximum value be achieved for $$ \frac{v_1}{v_0} = \frac{\gamma_1}{\gamma_1 + 1}, $$correspond value equal $$ \frac{2q}{c^2} = \frac{1}{\gamma_1^2 -1}, $$from which velocity wave $$ c = \sqrt{2q (\gamma_1^2 -1)} $$and pressure $$ p_1 = \frac{c^2}{v_0}\left( 1 - \frac{v_1}{v_0}\right) = \frac{2q (\gamma_1^2-1)}{v_0} \frac{1}{\gamma_1 + 1} = \frac{2q}{v_0} (\gamma_1 -1). $$


Second solution. The equation will have a unique solution if the dependence curves $p_1(v_1)$ touch each other. Let us equate the expressions for pressures and their derivatives with respect to $v_1$:


\begin{align*} \&\frac{c^2}{v_0} \left( 1 - \frac{v_1}{v_0}\right)= \frac{2q}{v_0} \frac{1}{\frac{\gamma_1 +1}{\gamma_1 - 1} \frac{v_1}{v_0} - 1},\\ \&\frac{c^2}{v_0^2} = \frac{2q}{v_0^2} \frac{\gamma_1 +1}{\gamma_1 -1}\frac{1}{\left(\frac{\gamma_1 +1}{\gamma_1 - 1} \frac{v_1}{v_0} - 1\right)^2}. \end{align*} Let's express the wave speed from the second equation and substitute it into the first: $$ c^2  = 2q \frac{\gamma_1 +1}{\gamma_1 -1}\frac{1}{\left(\frac{\gamma_1 +1}{\gamma_1 - 1} \frac{v_1}{v_0} - 1\right)^2}, $$$$ \frac{\gamma_1 +1}{\gamma_1 -1}\frac{1 - \frac{v_1}{v_0}}{\left(\frac{\gamma_1 +1}{\gamma_1 - 1} \frac{v_1}{v_0} - 1\right)^2} = \frac{1}{\frac{\gamma_1 +1}{\gamma_1 - 1} \frac{v_1}{v_0} - 1}, $$hence obtain equation on ratio $v_1/ v_0$: $$ \frac{\gamma_1 +1}{\gamma_1 -1}\left(1 - \frac{v_1}{v_0} \right) = \frac{\gamma_1 +1}{\gamma_1 - 1} \frac{v_1}{v_0} - 1, $$from which find $$ v_1 = v_0 \frac{\gamma_1}{\gamma_1 + 1}. $$Substitute this value in expression for $c^2$, obtain answer for $c$ and then $p_1$.

\par\smallskip\noindent\textbf{Answer:} $$ v_1 = \frac{\gamma_1}{\gamma_1 + 1} v_0, \quad p_1 =\frac{2q}{v_0} (\gamma_1 -1). $$
\end{solcard}

\begin{solcard}{B5}{Determine the speed of movement of the wave front $c$. Express your answer in terms of $q$ and $\gamma_1$.}{0.20}
\par\smallskip\noindent\textbf{Answer:} $$ c = \sqrt{2q (\gamma_1^2 -1)}. $$
\end{solcard}

\begin{solcard}{B6}{Calculate the speed of movement of the wave front during a detonation wave in a mixture of oxygen and hydrogen if, when it is heated to high temperatures, the following chemical reaction occurs: $$ 2\mathrm{H}_2+ \mathrm{O} _2 = \mathrm{H}_2 \mathrm{O} + \mathrm{H}^{+} + \mathrm{HO}^{-}. $$The actual process of hydrogen combustion is a combination of reaction, in which atoms of oxygen and hydrogen interact, therefore among the reaction products it is necessary to take into account the ion $\mathrm{H}^{+}$ and $\mathrm{HO}^{-}$ (this ion can be considered as a diatomic molecule). in mixture on 2 moles of hydrogen corresponds mole oxygen in correspondence with equation reaction. For this reaction consider/count, that $q \approx 15~\text{MJ}/\text{kg}$ for each kilogram of hydrogen consumed.}{0.60}

The reaction products are a mixture of monatomic, diatomic and polyatomic gases in equal proportions (by amount of substance), therefore the adiabatic index (indices 1,2,3 mark the values related to monatomic, diatomic and polyatomic gas, respectively) $$ \gamma_1 = \frac{C_{P1} + C_{P2} + C_{P3}}{C_{V1} + C_{V2} + C_{V3}} = \frac{10}{7}. $$Then velocity wave one can calculate by formula from previous part. At the same time needed take into account that the heat released is per unit mass of hydrogen. According to equation reaction, on 1 kg hydrogen corresponds 8 kg oxygen, therefore needed use value $q$ is 9 times less.

\par\smallskip\noindent\textbf{Answer:} $$ c = \sqrt{2 q (\gamma_1^2 - 1)} \approx 1.86 \cdot 10^{3}~\frac{\text{ m}}{\text{s}}. $$
\end{solcard}

\begin{solcard}{C1}{Obtain an equation for the light detonation adiabatic similar to point B2, that is, find the pressure $p_1$ in the plasma after the passage of the wave front. Express your answer in terms of $v_0$, $v_1$, $\gamma_1$, $I_0$.}{1.20}

Let's write down the law of conservation of energy from point $\mathrm{B2}$, replacing $q = I_0 v_0/c$ in it, expressing the speed of the wave $c$ through pressure and specific volumes: $$ I_0 v_0 \frac{\sqrt{1- v_1/v_0}}{\sqrt{p_1 v_0}}+ p_1 \frac{v_0^2}{2(v_0 - v_1)} = \frac{\gamma_1}{\gamma_1 -1}p_1 v_1 + \frac{v_1^2 }{v_0^2} p_1 \frac{v_0^2}{2(v_0 - v_1)} . $$Transform this equation $$ I_0 \frac{\sqrt{v_0 -v_1}}{\sqrt{p_1}} =\frac{ p_1 v_0}{2} \left( \frac{\gamma_1 +1}{\gamma_1 - 1} \frac{v_1}{v_0} - 1\right) , $$ from where we get

\par\smallskip\noindent\textbf{Answer:} $$ p_1 =\left( \frac{2I_0 \sqrt{v_0 -v_1}}{\frac{\gamma_1 +1}{\gamma_1 - 1} v_1 - v_0} \right)^{2/3}. $$
\end{solcard}

\begin{solcard}{C2}{Determine the parameters ($p_1$ and $v_1$) of the Champin-Jouguet point for the light-detonation adiabat. Express your answer using $v_0$, $\gamma_1$, $I_0$.}{1.40}

Let us again write down the equality of the pressures obtained from the formula for the light-detonation adiabat and from the speed of the wave: $$ \left( \frac{2I_0 \sqrt{v_0 -v_1}}{\frac{\gamma_1 +1}{\gamma_1 - 1} v_1 - v_0} \right)^{2/3} = \frac{c^2}{v_0} \left(1  - \frac{v_1}{v_0} \right). $$Transform this equation: $$ \left(1  - \frac{v_1}{v_0} \right)^{2/3} \left( \frac{\gamma_1 +1}{\gamma_1 - 1} v_1 - v_0 \right)^{2/3} =\frac{ (2 I_0 v_0)^{2/3} }{c^2}. $$Raising the equation to the power $3/2$, we again obtain a quadratic equation for $v_1/v_0$: $$ \left(1- \frac{v_1}{v_0} \right)\left( \frac{\gamma_1 +1}{\gamma_1 - 1} \frac{v_1}{v_0} - 1\right) = \frac{2I_0 v_0}{c^3}. $$ The left side is the same as in the case of a shock wave, therefore the maximum is reached at $$ v_1 = \frac{\gamma_1}{\gamma_1 + 1} v_0, $$ and the wave velocity can be found from the relation $$ \frac{2 I_0 v_0}{c^3} = \frac{1}{\gamma_1^2 -1}, \quad c = \left(2 I_0 v_0 (\gamma_1^2 -1) \right)^{1/3}. $$Then pressure $$ p_1 = \frac{\left(2 I_0 v_0 (\gamma_1^2 -1) \right)^{2/3}}{v_0} \frac{1}{\gamma_1+ 1} = \left( \frac{4 I_0^2}{v_0}\frac{(\gamma_1 -1)^2}{\gamma_1 + 1}\right)^{1/3}. $$For comparison we also plot on one graph the light-detonation adiabat and the dependences $p_1(v_1)$ at constant speed of sound.

\par\smallskip\noindent\textbf{Answer:} $$ v_1 = \frac{\gamma_1}{\gamma_1 + 1} v_0, \; p_1 =  \left( \frac{4 I_0^2}{v_0}\frac{(\gamma_1 -1)^2}{\gamma_1 + 1}\right)^{1/3}. $$
\begin{figure}[H]
  \centering
  \includegraphics[width=0.72\textwidth]{T4_S_EN_assets/fig2.jpg}
\end{figure}
\end{solcard}

\begin{solcard}{C3}{Determine the speed of movement of the wave front $c$ of the light detonation wave. Express your answer in terms of $\gamma_1$, $I_0$, $\rho_0$. Calculate its numerical value.}{0.40}

Let's use the formula for the wave front speed from the previous paragraph, substitute the intensity value $I_0 = P/\pi r^2 \sim 10^{15}~\text{W}/\text{m}^2$ and the value $v_0 = 1/\rho_0$ into it.

\par\smallskip\noindent\textbf{Answer:} $$ c =  \left(2 I_0 v_0 (\gamma_1^2 -1) \right)^{1/3} \approx 100 ~\frac{\text{km}}{\text{s}} $$
\end{solcard}

\end{document}